\begin{afterword}
\summaryhead

The post continues from the author's previous one, which suggested that humans may have evolved to seek power believing it is for the common good, and then to be corrupted by it. If so, the feeling that seizing power would be altruistic is weak evidence that it would be. The human remedy is a rule: do not cheat to seize power, and do not murder, even for the good of the tribe.

The author then takes up the philosopher's case in which pushing one person in front of a train is certain to save five. One reply is that a human cannot know such facts, so only an AI free of the bias could rightly push. The author calls this a dodge, since the world can force such choices, but holds that humans should keep their prohibitions while a well-calibrated Friendly AI might push, because the tendency to be corrupted by power is a biological adaptation an AI would not share. Minds biased the other way would need a reversed rule. ``The end does not justify the means'' is therefore ``reflective consequentialism''.

\responsehead

The post's central claim is sound, and it is old. A person who knows that their judgment is biased in their own favour has reason to keep rules that forbid what the bias recommends, even at some cost in cases where the bias is not at work. The post frames this as a problem outside classical decision theory, which is fair. In ethics it is the standard indirect consequentialist position. Mill answered in 1861 the objection that a utilitarian ``when under temptation, will see an utility in the breach of a rule''. Sidgwick suggested that calculation might be left to an elite better at it. Hare's archangel, with superhuman knowledge and no weaknesses, needs none of the rules ordinary people must keep, and the post's AI takes the archangel's place. What the post adds is an evolutionary account of why the human bias runs toward power, and the application to AI.

That account is the weakest step. The post's first paragraph says humans ``may have evolved'' the tendency, and the previous day's post offered it with ``Call it a just-so story if you must'' and raised an objection from egalitarian hunter-gatherer bands. By the Friendly AI paragraph it has become ``a specific biological adaptation, supported by specific cognitive circuits''. The case for exempting an AI from the rule depends on that sentence, and no evidence for it is given. Nor does the post give the ``historical statistics'' said to justify the human prohibitions.

The trolley passage rejects less than it seems to. The post calls a reply a dodge, then reaches the same verdicts on the same ground: humans keep the prohibition because their judgment cannot be trusted, and an AI without the bias may push. What it rejects, as far as the text shows, is the reply's refusal to consider such cases. It grants that the world can force them, and notes that legal systems already weigh innocent against guilty.

The post takes a side in an old debate. The main published objection to two-level views is that a person who knows the rules are only instruments may not hold them firmly, or may not keep the two kinds of thinking apart. The next post in the book, ``Ethical Injunctions'', takes up a version of that worry.

In short: a sound and long-established idea, indirect consequentialism, given an evolutionary rationale that the post states as fact though its own opening and the previous day's post hedge it, and on which its exemption for AI depends.
\end{afterword}
